\begin{enumerate}
    \item 本研究以 He-Ne 雷射、TSL2591 光感測器與樂高旋轉臂建立低成本多角度散射量測系統；同一樣品兩輪量測所得 $C$ 的差異多在 $4\%$ 以內。
    \item 在 $\tau=0.5$ 至 $8$ 的範圍內，$C=I(30^\circ)/I(125^\circ)$ 隨光學厚度增加而下降；相同 $\tau$ 下，$2.16\,\mu\mathrm{m}$ 微球的 $C$ 約為 $0.492\,\mu\mathrm{m}$ 微球的 $1.6$ 至 $2.4$ 倍，顯示兩角強度比在單次散射至多重散射的過渡範圍仍保有粒徑資訊。
    \item $2.16\,\mu\mathrm{m}$ 微球的 $C$ 下降較慢，方向上符合較高 $g$ 所代表的較長傳輸平均自由徑；但下降速率差距僅約 $10\%$ 至 $15\%$，顯示 $C$ 仍受完整散射角度分布影響。
    \item 蒙地卡羅模擬重現了 $C$ 隨 $\tau$ 下降及粒徑差異的趨勢。當 $\tau\leq1$ 時，實驗與模擬比值為 $0.83$ 至 $1.05$；當 $\tau\geq2$ 時，實驗值高於模擬約 $10\%$ 至 $33\%$。前向角附近線材反射是可能原因之一，仍待驗證。
\end{enumerate}
